% =============================================================================
% Project : THE_PHILOSOPHY_OF_VIABILITY_LIBRARY_PROJECT
% Book    : Book I — Ontology of Viability
% Book ID : VS.B01
% Part    : Part I — Identity and Manifestation
% Chapter : Chapter 08 — Comparative Philosophical Examination
% Chapter ID : VS.B01.P01.C08
%
% File :
%   4_ASSETS/illustrations/VS.B01.P01.C08.S05.SS08.P005.IL01.matter_form_and_artifact.tex
%
% Section / Chapter / Asset:
%   Matter, Form, and Artifact
%
% Structural Identifier :
%   VS.B01.P01.C08.S05.SS08.P005.IL01
%
% Asset Type : IL
%
% Purpose :
%   To provide an aesthetically composed visual analogy for the
%   distinction between material substrate, organization, and
%   the resulting artifact.
%
% Dependencies :
%   TikZ; positioning library, loaded by the project preamble.
%
% Representational limit :
%   The illustration concerns an artifact and must not be taken
%   as a complete representation of natural substances or as
%   evidence for hylomorphism.
%
% =============================================================================

\begin{figure}[htbp]
\centering

\begin{tikzpicture}[
    material/.style={
        draw=ViabilityGold,
        fill=ViabilityPaleGold,
        rounded corners=2mm,
        line width=0.8pt,
        minimum width=34mm,
        minimum height=17mm,
        align=center,
        text=ViabilityBlack
    },
    form/.style={
        draw=ViabilityDarkBlue,
        fill=ViabilityPaleBlue,
        rounded corners=2mm,
        line width=0.8pt,
        minimum width=34mm,
        minimum height=17mm,
        align=center,
        text=ViabilityBlack
    },
    result/.style={
        draw=ViabilityGreen,
        fill=ViabilityPaleGreen,
        rounded corners=2mm,
        line width=0.8pt,
        minimum width=43mm,
        minimum height=21mm,
        align=center,
        text=ViabilityBlack
    },
    flow/.style={
        -{Stealth[length=2.5mm]},
        line width=1pt,
        draw=ViabilityGray
    }
]

\node[material] (wood) {
    \textbf{Material}\\
    Wood
};

\node[form, right=17mm of wood] (organization) {
    \textbf{Organization}\\
    Shape and arrangement
};

\node[result, below=18mm of $(wood)!0.5!(organization)$] (artifact) {
    \textbf{Artifact}\\
    A constituted wooden object
};

\draw[flow] (wood.south) -- ++(0,-6mm) -| (artifact.west);
\draw[flow] (organization.south) -- ++(0,-6mm) -| (artifact.east);

\node[
    below=7mm of artifact,
    text=ViabilityGray,
    align=center,
    font=\small,
    text width=105mm
] {
    A limited analogy for the relation between what a thing is made
    from and the organization through which it constitutes an artifact.
};

\end{tikzpicture}

\caption{A visual analogy for the distinction between material
substrate and organization in an artifact. The illustration is
not a complete model of Aristotle's account of natural substances.}
\label{fig:aristotle-matter-form-artifact}

\end{figure}

% =============================================================================
% END OF ILLUSTRATION IL01
% =============================================================================